\section{The prescribed-initial proper iteration proof}
\label{annex:p06-proper-iteration-lemma}

\begin{annexitem}{annex:p06.pil.lifting}{thm:p06-pil-lifting}{Small names and the actual lifted model}
\begin{proof}
First record the hereditary-name interface at the declared height.
Since \(|P|\leq\kappa<\chi\), possible-value antichains show that
\(P\) preserves regularity of \(\chi\): fewer than \(\chi\) coordinates
each have at most \(\kappa\) possible values, so their union is bounded
below the ground regular \(\chi\). A name in \(H_\chi\) evaluates
to an element of \(H_\chi^{V[G]}\), since its transitive closure is
covered by the evaluations of its fewer-than-\(\chi\) member names.

Conversely code a hereditary-small extension set by a well-founded
extensional relation on some ordinal \(\mu<\chi\), with a distinguished
root. Choose a nice name for that relation: at each pair in
\(\mu\times\mu\), a maximal deciding antichain has size at most
\(\kappa\). Decide ground bounds on \(\mu\) and on all relation ranks
below \(\chi\); the antichain possible-value bound makes a common
\(\rho<\chi\) available below the condition being considered.
Mix with a trivial relation outside it. Construct collapse-value names
\(v_{i,\eta}\), \(i<\mu,\eta<\rho\), in increasing \(\eta\): include
\((v_{j,\zeta},a)\) when \(\zeta<\eta\) and \(a\) forces
\(jRi\), rank of \(j\) equal to \(\zeta\), and rank of \(i\)
equal to \(\eta\). Induction on \(\eta\) gives the actual collapse
value at every active rank. Mix the distinguished root's rank values.
There are fewer than \(\chi\) indices and pairs throughout, so the
resulting name is hereditarily smaller than \(\chi\). This proves
the needed small-name representation, not merely a rank-cutoff assertion.

Fixed-formula forcing restricted to names in \(H_\chi\) is definable
there and has the truth/density/witness clauses for \(H_\chi^{V[G]}\).
For the atomic recursion, the transitive closures of the finitely many
input names and \(\mathcal P(P)\) have size below \(\chi\); the latter
is available because \(\chi>2^\kappa\). Its local recursion tables
therefore belong to \(H_\chi\). Formula induction uses the small-name
representation for existential witnesses. Collection is valid in
\(H_\chi\): fewer than \(\chi\) witnesses, each hereditary-small,
have hereditary-small union by regularity. Thus antichain witness
selection stays there; no Power Set axiom for the whole \(H_\chi\)
or set of every name at an unbounded rank is invoked.

Now \(N[G]\) is countable. For names \(\bar\tau\in N\) and a fixed
formula, the set of conditions deciding the existential statement and,
on its positive side, naming a hereditary-small witness, is dense and
definable in \(H_\chi\), hence belongs to \(N\). A generic containing
a master meets its intersection with \(N\). If the existential is true,
the deciding condition cannot take the negative side; elementarity selects
its witness name in \(N\). Its evaluation lies in \(N[G]\). This is
the full Tarski--Vaught check for \(N[G]\prec H_\chi^{V[G]}\).

For a ground \(a\in N\), use the dense set deciding either
\(\tau\notin\check a\) or \(\tau=\check x\) for some \(x\in a\),
with \(\tau\in N\). A master generic meets it in \(N\); when the
actual value is in \(a\), the selected ground value is in \(N\) by
elementarity and uniqueness. Thus \(N[G]\cap a=N\cap a\).
For the full ground trace, let \(S_\tau\) be all ground values which
some condition forces equal to \(\tau\). There are at most \(|P|\)
such values, one per deciding condition. Their hereditary sizes are
uniformly bounded below \(\chi\): every member of their transitive
closures is a possible ground value of a member name in the transitive
closure of \(\tau\), giving at most
\(|\operatorname{tc}(\tau)|\cdot|P|\) values. Thus \(S_\tau\in H_\chi\),
and its atomic-forcing definition makes it an element of \(N\).
If \(\tau^G\) is ground, truth puts it in \(S_\tau\), hence in \(N\)
by the ground-set trace. Conversely check names put every member of
\(N\) in \(N[G]\). Ordinals are ground, giving the ordinal trace.
All of these arguments have dense-set formulations below the master,
so supply the corresponding forced assertions used at successor stages.
\end{proof}
\end{annexitem}

\begin{annexitem}{annex:p06.pil.tail-representation}{thm:p06-pil-tail-representation}{The quotient, not an invented fresh support iteration}
\begin{proof}
The canonical quotient condition name has \(\check u\) supported by
\(u\mathbin{\upharpoonright}\alpha\). Its evaluation is exactly the
displayed ground-represented poset. The map from \(P_\beta\) into
\(P_\alpha*\dot T\) sends \(u\) to its restriction and the name
equal to \(\check u\) below that restriction, mixed with the unit
elsewhere. It preserves and reflects order and incompatibility. For
density, decide a proposed quotient condition as a ground \(u\), refine
its forced generic-prefix membership to a common prefix below its
restriction, and apply the exact initial-segment lifting property.
The lifted full condition maps below the proposed two-step pair. This
proves factorization using the supplied two-step construction.

The relation \(\leq_s\) is transitive: below an alleged first condition,
first find a common strengthening with the second, then one with the
third. Mutual comparison quotients it to a separative partial order.
Compatibility and dense requirements are unchanged. Conditions with
the same tail and prefixes in \(G\) are equivalent: below either,
choose a common generic prefix with the other and paste it onto the
stronger tail. This gives a common old-order strengthening in the
quotient. No closure of an iterand is needed for this pasting.

The canonical quotient and class map are elements of \(N[G]\), since
their names are definable from the initial-segment system in \(N\)
and have hereditary size at most \(2^\kappa<\chi\).
Elementarity selects a representative \(u\in T(G)\cap N[G]\)
of the named model class. It is a ground \(P_\beta\)-condition,
so the ground trace makes it an element of \(N\). The chosen raw
section is in the full iteration datum, hence gives a raw representative
in \(R_\beta\cap N\). Its countable support belongs to \(N\),
and the countable-element lemma makes that support a subset of \(N\).
Antichain mixing chooses the variable raw name uniformly below \(p\);
its values, not necessarily the name itself, belong to \(N\).
\end{proof}
\end{annexitem}

\begin{annexitem}{annex:p06.pil.successor-master}{thm:p06-pil-successor-master}{Select the master name and verify all dense sets}
\begin{proof}
Lifting makes \(N[G]\) a countable elementary submodel at an admissible
height. Properness of the actual iterand gives a master below the
actual \(\dot a^G\in N[G]\). The forcing assertion is uniform below
\(p\): existential witnesses can be selected on a maximal antichain,
bounded by Collection and mixed. The bounded condition-name lemma
then gives \(\dot b\) with \(p\) forcing both masterhood and the
required inequality. There is no request that this selected name be in \(N\).

For a dense open two-step \(D\in N\), its right-hand dense set in
the first extension belongs to \(N[G]\), by the actual two-step
definition. A master \(\dot b^G\) makes its intersection with
\(N[G]\) predense. Elementarity of \(N[G]\) chooses a pair witnessing
membership in that right-hand dense set; its ground coordinates and
name lie in \(N\) by the ground trace. Thus this pair belongs to
\(D\cap N\). To make the predensity verification nonvacuous, start
with any stronger two-step condition. In the first forcing select a
common right strengthening with such a model witness, decide its
ground witness pair, refine the first coordinate below its forced
generic support, and name the common right condition. This is a ground
common two-step strengthening with a member of \(D\cap N\).
Therefore \((p,\dot b)\) is a master, with no generic-existence
assumption replacing predensity.
\end{proof}
\end{annexitem}

\begin{annexitem}{annex:p06.pil.extension}{thm:p06-pil-extension}{Support-controlled successor and limit induction}
\begin{proof}
Use the quotient lemma to select a raw name \(\dot\rho\) valued in
\(R_\beta\cap N\), with its initial class forced into the generic.
Mix with the raw unit outside the starting condition. For \(\xi\geq\alpha\)
there is an explicit \(P_\xi\)-name \(\widehat\rho_\xi\) for the
value of the selected raw coefficient: take an antichain in \(P_\alpha\)
deciding \(\dot\rho=u\), lift that antichain to \(P_\xi\), and mix
the names \(u(\xi)\) on its supports. Normalize into the bounded
coordinate name set. This name evaluates to \(u(\xi)^{G_\xi}\)
when \(u=\dot\rho^{G_\alpha}\). At indices outside \(N\) it is the
unit, because every raw value's support is contained in \(N\).

We prove by induction on the endpoint a stronger statement: the constructed
raw \(r\), after \(\alpha\), has support in \(N\), and
\[
 [r\mathbin{\upharpoonright}\xi]\forces
       r(\xi)\leq\widehat\rho_\xi\qquad(\alpha\leq\xi<\beta).
 \tag{*}
\]
This implies the desired quotient comparison. Indeed, in an initial
generic containing \(p\), let \(u\) be the selected raw value.
Below any old quotient strengthening \(v\) of \(r\), choose a generic
prefix below its prefix and \([u\mathbin{\upharpoonright}\alpha]\),
also deciding \(\dot\rho=u\). Paste that prefix onto the tail of
\(v\). Monotonicity of (*) makes it a full-stage common strengthening
of \(v\) and \(u\). Thus the separative class of \(r\) is below that
of \(u\). The same argument has a ground dense formulation: refine a
prefix to decide the raw value, refine its forced generic membership
to a common prefix, and paste the stronger tail. It proves
\(r\forces[\dot\rho]\in\dot G_\beta\), not just a vacuous statement
about possible generics. At any smaller stage it likewise gives forced
membership of the corresponding raw initial class.

At equal endpoints no tail is changed. For \(\beta=\eta+1\),
\(\eta\in N\). Apply the induction hypothesis to the raw restriction
of \(\dot\rho\) to \(\eta\), preserving the prescribed initial
condition. In a generic containing the resulting master \(r_\eta\),
the raw value is in \(N\), its coefficient at \(\eta\) is a name
in \(N\), and its evaluation belongs to \(N[G_\eta]\).
The successor-master lemma selects an actual bounded name for an
iterand master below \(\widehat\rho_\eta\). Append it. The two-step
master test proves masterhood, (*) holds at the new coordinate, and
the only new support coordinate is \(\eta\in N\). Names may be outside
\(N\); their actual model values, which are what properness uses, are
inside \(N[G_\eta]\).

Now let \(\beta\) be limit and put \(\delta=\sup(N\cap\beta)\).
Choose increasing \(\alpha_n\in N\cap\beta\), \(\alpha_0=\alpha\),
cofinal in \(\delta\), and enumerate all raw dense open sets
\(D_n\in N\) of \(R_\beta\). Raw pullbacks of quotient dense sets
belong to this list, so this suffices for the master test. Construct
masters \(p_n\in P_{\alpha_n}\) and raw-condition names \(\rho_n\)
with values in \(R_\beta\cap N\), their initial classes forced into
\(G_{\alpha_n}\). Start with the prescribed \(p,\rho\).

Given \(p_n,\rho_n\), one can choose, in the first extension, a raw
\(u\in D_n\cap N\) stronger than \(\rho_n\), whose restriction to
\(\alpha_n\) is in the generic. Here is the exact dense-set argument.
For each fixed raw value \(w\in N\), project the part of \(D_n\)
below \(w\) to \(P_{\alpha_n}\). The exact restriction-lifting
property makes this projection dense below \([w\mathbin{\upharpoonright}
\alpha_n]\). Add the incompatible conditions to make a global dense
set; it belongs to \(N\). The initial master meets its intersection
with \(N\), on the compatible side. Elementarity supplies a raw
witness in \(D_n\cap N\) with that generic prefix. Mixing selects a
name \(\sigma_n\) for this \(u\), uniformly below \(p_n\).

Apply the induction hypothesis at endpoint \(\alpha_{n+1}<\beta\)
to \(\sigma_n\mathbin{\upharpoonright}\alpha_{n+1}\), with prescribed
initial \(p_n\). It gives a master \(p_{n+1}\) whose restriction is
exactly \(p_n\), satisfying (*) for that block. Let \(\rho_{n+1}\)
be the lifted name \(\sigma_n\) at this new stage. It has values in
\(D_n\cap N\), is forced stronger than \(\rho_n\), and its initial
class is forced into \(G_{\alpha_{n+1}}\). Choose raw representatives
of \(p_{n+1}\) with literally the fixed preceding prefix; replacing an
equivalent prefix changes no quotient condition.

Take the union of these coherent raw initial functions through
\(\delta\), and put unit coordinates on \([\delta,\beta)\).
Its support is contained in
\(\operatorname{supp}(p)\cup(N\cap[\alpha,\delta))\), a countable
ground set. Every restriction below \(\delta\) is a restriction of
some \(p_n\); hence the union is a legitimate countable-support limit
condition. If \(\delta<\beta\), extension by units is legitimate as
well. The original support of \(p\), even if not in \(N\), is entirely
below \(\alpha\); it has never been discarded.

We verify domination rather than assume it from coherence. For a
coordinate in block \([\alpha_n,\alpha_{n+1})\), the block construction
forces its new coefficient below the selected \(\rho_{n+1}\) coefficient.
The selected raw conditions are forced descending. Their initial
classes at the relevant smaller stages are forced generic members;
refining below such a member is dense, so its raw coordinate inequalities
transfer to the constructed prefix by monotonicity and density. Thus
the coefficient is below every earlier \(\rho_j\) coefficient, after
interpreting their earlier names through the coherent projections.
This proves (*) for the original \(\rho\), and, on the remaining tail,
for each later \(\rho_j\). At an index outside \(N\) both relevant
coefficients are units. At or above \(\delta\) every selected raw model
condition has unit coordinates, because its countable support is a
subset of \(N\cap\beta\). Hence no unhandled support survives beyond
the trace, even when \(\delta<\beta\).

The dense common-prefix/pasted-tail argument from (*) now proves
\(r\forces[\rho_{n+1}]\in\dot G_\beta\) for every \(n\).
Its value lies in \(D_n\cap N\). To recover ground predensity, below
any strengthening of \(r\) refine the bounded name to a fixed member
of \(D_n\cap N\); canonical generic membership then supplies a common
strengthening with that fixed member. Thus \(D_n\cap N\) is predense
below \(r\). This proves masterhood and the quotient comparison,
completing the induction without any arbitrary proper lower bound.
\end{proof}
\end{annexitem}

\begin{annexitem}{annex:p06.pil.model-masters}{thm:p06-pil-model-masters}{The trivial initial stage produces the stated masters}
\begin{proof}
At \(\alpha=0\), the unit is a master. A raw representative of
\(u\in P_\beta\cap N\) can be selected in \(N\) by the fixed
section in the iteration datum; use its constant name. The extension
lemma gives \(r\) master, and (*) at every coordinate makes it a
literal raw strengthening of that representative. Its quotient condition
therefore strengthens \(u\). This proves exactly the stated
iteration-datum model test, with its parameters retained.
\end{proof}
\end{annexitem}
